\documentclass[UTF8,11pt,a4paper]{ctexart}
\usepackage{../../tex/ai_infra_course}
\setmonofont{Noto Sans Mono CJK SC}

\title{\Huge RMSNorm：\[.35em]\Large 起一个稳压器的作用}

\begin{document}
\maketitle
\section{摘要}
RMSNorm 的主要作用是稳定 Transformer 中各层输入向量的整体尺度。可以把它类比成电路里的“稳压器”：前面传来的数值可能忽大忽小，而 RMSNorm 会在共同归一化后由逐维权重再调尺度，从而让模型各层的计算更加稳定。
\section{文本到向量的变化}
先复习下文本是怎么变成向量的，懂的可以直接跳

先看一个中文句子：
\begin{center}
\texttt{乒乓球拍卖完了。}
\end{center}

如果我们要通过某种方法把这段文字转换成可计算的数学形式，第一步要把它切成一个个小片段，这些小片段叫 \textbf{token}。切分之后每个 token 会到一个词表里查出对应的整数编号。

例如，为了方便理解，我们暂时假设 Tokenizer 把它切成：
\[
[\texttt{乒乓球},\texttt{拍卖},\texttt{完},
 \texttt{了},\texttt{。}]
\longrightarrow
[18432,9271,643,352,17].
\]

<details>
<summary>Tokenizer、词表与 token ID 是怎么来的？</summary>

原始输入首先只是一个字符串，例如：

\[
\texttt{"乒乓球拍卖完了。"}
\]

它会先被送入 \textbf{Tokenizer（分词器）}。Tokenizer 的作用是按照一套固定规则，把字符串切分成若干个 \textbf{token}。

中文在这里有一个很直观的问题：原始字符串本身并没有空格明确告诉我们“词”在哪里结束。

例如：

\[
\texttt{乒乓球拍卖完了}
\]

从人类语言的角度，就至少可以出现两种不同的划分：

\[
[\texttt{乒乓球拍},\texttt{卖},\texttt{完},\texttt{了}]
\]

可以理解成“乒乓球拍已经卖完了”；而

\[
[\texttt{乒乓球},\texttt{拍卖},\texttt{完},\texttt{了}]
\]

又可以理解成“乒乓球的拍卖结束了”。

也就是说，同一串连续的汉字并不存在天然写在字符串里的唯一“词语边界”。

不过，真实语言模型的 Tokenizer 并不是先像人一样理解句子含义，再决定应该按照哪个中文词语切分。它会根据训练语料统计出的词表和合并规则处理字符序列：某种连续片段出现得越稳定，就越可能作为一个 token；同一串汉字也可能因为词表不同而得到不同切分。因此，真实的 token 边界不一定等于我们通常理解的中文分词。
因此，真实的 token 边界也不一定等于我们通常理解的“中文分词”。

对于英语，最简单的方法就是按照空格划分，但是实际上Tokenizer不关心字符是空格、标点还是具体字母和字，所以“I love you.”,中可能出现“love”的token，也可能出现“ love”的token（含前导空格）。甚至完整单词，比如“unchanged”，也可能因为“changed”的单独出现频率太高而被分为“un”，“changed”。
所以，token实际上是 Tokenizer 根据自己的词表和切分规则得到的基本符号单位。

Tokenizer 同时拥有一个固定的 \textbf{词表（vocabulary）}。可以把词表简单理解成一张

\[
\text{token}\longleftrightarrow\text{整数编号}
\]

的映射表，例如：

\[
\texttt{乒乓球}\rightarrow18432,\qquad
\texttt{拍卖}\rightarrow9271,\qquad
\texttt{。}\rightarrow17.
\]

因此，分词完成后，每个 token 都会根据词表被替换成对应的整数编号，即 \textbf{token ID}：

\[
[\texttt{乒乓球},\texttt{拍卖},\texttt{完},
\texttt{了},\texttt{。}]
\longrightarrow
[18432,9271,643,352,17].
\]

这里的编号只是为了说明数据流而举出的例子，真实编号由具体 Tokenizer 的词表决定。

这些整数本身没有大小或语义上的数值关系，例如编号 \(18432\) 并不意味着它比编号 \(9271\)“更大”或“更重要”，它们本质上只是为了方便寻址的一个索引。

接下来，这些 token ID 会被送入 \textbf{Embedding 表}。Embedding 表可以看作一个很大的矩阵，每个 token ID 对应其中的一行：

\[
\text{token ID}
\longrightarrow
\text{Embedding 表中的对应行}
\longrightarrow
\text{浮点向量}.
\]

因此，从原始文字到模型真正参与计算的数据，完整的数据流是：

\[
\boxed{
\text{原始字符串}
\rightarrow
\text{Tokenizer 分词}
\rightarrow
\text{tokens}
\rightarrow
\text{词表查询}
\rightarrow
\text{token IDs}
\rightarrow
\text{Embedding 查询}
\rightarrow
\text{浮点向量}
}
\]
也就是说，Tokenizer 负责把“文字”变成“整数索引”，Embedding 再负责把“整数索引”变成神经网络真正能够计算的浮点向量。

</details>

现在模型手里是一个$1\times5$的向量：1 个句子，5 个 token。每个 token 有对应的编号，再从 Embedding 表里取出对应编号索引的一个向量。为了说明方便，我们这里把向量长度设成 4，那么这句话的 5 个 token 从 Embedding 表中取出的向量可能长得类似于这样：

\begin{array}{c|c|c}
\text{token} & \text{token ID} & \text{Embedding 向量} \\
\hline
\texttt{乒乓球} & 18432 & [\,0.10,\;0.25,\;-0.05,\;0.20\,] \\
\texttt{拍卖}   & 9271  & [\,-0.30,\;0.15,\;0.40,\;0.05\,] \\
\texttt{完}     & 643   & [\,0.12,\;-0.18,\;0.22,\;0.31\,] \\
\texttt{了}     & 352   & [\,0.08,\;0.11,\;-0.14,\;0.27\,] \\
\texttt{。}     & 17    & [\,-0.06,\;0.09,\;0.03,\;-0.12\,]
\end{array}

把这 5 个向量按 token 的顺序堆起来，就得到了这个句子的 Embedding 输出：

\begin{bmatrix}
 0.10 &  0.25 & -0.05 &  0.20 \\
-0.30 &  0.15 &  0.40 &  0.05 \\
 0.12 & -0.18 &  0.22 &  0.31 \\
 0.08 &  0.11 & -0.14 &  0.27 \\
-0.06 &  0.09 &  0.03 & -0.12
\end{bmatrix}


如果再把 batch 维度算上，那么它的形状就是
\[
(B,T,D)=(1,5,4).
\]
这里 1 表示一个句子，5 表示 5 个 token，4 表示每个 token 用一个 4 维向量表示。




\section{尺度漂移}
经过 Embedding 之后，每个 token 都已经变成了一个浮点向量。接下来，这些向量会被送入 Transformer 的一层层计算中，并在计算过程中不断产生新的中间向量。

随着层数加深，这些向量的数值大小可能发生明显变化。例如某一层中，一个 token 的向量可能是
\[
x=[1,2,3,4],
\]
但另一次训练，它可能变成
\[
x'=[10,20,30,40].
\]

虽然两个向量各个维度之间的比例完全相同，但整体尺度相差了 10 倍。这样的尺度变化如果不断累积，会让后续计算面对很不稳定的数值范围。

所以在向量进入后续计算之前，必须先对它的整体尺度进行控制，使不同层、不同 token 的向量不会因为整体数值过大或过小而给后续计算带来额外困难。

注意力层里还要算\textbf{点积}（两个向量逐项相乘再求和）和\textbf{softmax}（把一组数变成概率）；前馈层里有\textbf{激活函数}（对每个数单独做非线性映射）。如果展示不知道这些什么意思也没关系，只需要明白这些运算对数值大小都很敏感。如果很多层连续把不同尺度的向量往后传，同一个子层一会儿接收 $0.1$ 量级的数，一会儿接收 $100$ 量级的数，训练会变得不可控，低精度计算的风险也会变大。

所以我们需要把每个向量的整体大小控制住，同时保留内部各维度的相对关系。也就是归一化干的事情。

\section{LayerNorm}

LayerNorm 是较早用在 Transformer 里的归一化方法。它对每个 token 的向量分别算两个统计量：均值（平均数）和方差（数值偏离均值的程度）。

还是用 $x=[1,2,3,4]$ 当例子。四个数的均值是
\[
\mu=\frac{1+2+3+4}{4}=2.5.
\]
每个数减掉均值：
\[
x-\mu=[-1.5,-0.5,0.5,1.5].
\]
这四个数的方差是
\[
\sigma^2=\frac{\sum_{i=1}^{D}{(x_i-\mu)^2}}{D}=\frac{(-1.5)^2+(-0.5)^2+0.5^2+1.5^2}{4}=1.25，
\sigma\approx1.118.
\]
然后再让刚才得到的每个数都除以这个标准差$\sigma$：
\[
\frac{x-\mu}{\sigma}
=
\left[
\frac{-1.5}{1.118},
\frac{-0.5}{1.118},
\frac{0.5}{1.118},
\frac{1.5}{1.118}
\right].
\]

得到：
\[
\frac{x-\mu}{\sigma}
\approx
[-1.342,-0.447,0.447,1.342].
\]
现在这个新向量有两个特点：

\[
\text{均值$\mu$}=0,
\qquad
\text{方差$\sigma^2$}=1.
\]

通过这套步骤

\[
\boxed{
x
\longrightarrow
x-\mu
\longrightarrow
\frac{x-\mu}{\sigma}
}
\]

第一步减去均值，负责把整个向量“移到 0 附近”；第二步除以标准差，负责把向量的整体尺度调整到一个固定范围。
这实际上就是标准化，把一组均值、尺度各不相同的数字，转换成一组以 0 为中心、具有统一尺度的数字。

上面其实就是LayerNorm的核心思想了。而LayerNorm完整定义还会加入 $\gamma/\beta$ 两个可训练的缩放/平移参数，并用 $\varepsilon$ 防止除零和数值不稳定。

\[
\boxed{
y_i
=
\gamma_i
\frac{x_i-\mu}
{\sqrt{\sigma^2+\varepsilon}}
+
\beta_i
}
\]

LayerNorm 完整的数据流是

\[
\boxed{
x_i
\longrightarrow
x_i-\mu
\longrightarrow
\frac{x_i-\mu}{\sqrt{\sigma^2+\varepsilon}}
\longrightarrow
\gamma_i
\frac{x_i-\mu}{\sqrt{\sigma^2+\varepsilon}}
+\beta_i
}
\]

其中：

\begin{itemize}[leftmargin=2em]
    \item $\varepsilon$ 是一个很小的正数，用来避免分母为 0 或过小时出现数值不稳定，常见形式是 $\sqrt{\sigma^2+\varepsilon}$；
    \item $\gamma_i$ 是第 $i$ 个维度的可训练缩放参数；
    \item $\beta_i$ 是第 $i$ 个维度的可训练平移参数。
\end{itemize}

其中 $\mu$、$\sigma^2$ 都是根据当前 token 的向量现场计算出来的，而 $\gamma$ 和 $\beta$ 则是模型参数，会随着训练不断更新。

例如，如果向量长度为 4，那么 LayerNorm 自己还保存着两组长度为 4 的参数：

\[
\gamma=[\gamma_1,\gamma_2,\gamma_3,\gamma_4],
\qquad
\beta=[\beta_1,\beta_2,\beta_3,\beta_4].
\]

在 LayerNorm 刚初始化时，通常有

\[
\gamma=[1,1,1,1],
\qquad
\beta=[0,0,0,0].
\]

因此在训练刚开始时，LayerNorm 的输出基本就是前面算出的标准化结果；之后模型可以通过训练 $\gamma$ 和 $\beta$，重新调整每一个维度的尺度和位置。
\section{RMSNorm ：只缩放，不居中}

RMSNorm 的想法更简单：Transformer 子层最需要的是稳定的输入尺度，减均值不一定是必需的。所以它不再强制均值为 0，只保留缩放这一步。

不减均值以后，原来围绕均值算的标准差就不适用了(因为不仅分子中的 x−μ 被去掉了，标准差 σ 本身也是通过 x−μ 计算出来的)。RMSNorm 改用\textbf{均方根（root mean square，RMS）}来衡量尺度：
\[
\operatorname{RMS}(x)
=\sqrt{\frac{1}{D}\sum_{i=1}^{D}x_i^2}.
\]
平方让正负数都变成正贡献，而取平均得到每个分量平均平方大小，开平方把单位恢复到原来的大小。

直观地理解，也可以认为 RMS 和标准差的形式非常相似，区别主要在于计算“距离”时选择的参考点不同：

\begin{itemize}[leftmargin=2em]
    \item $\sigma$：以均值 $\mu$ 为参考点，衡量各个分量距离 $\mu$ 的整体大小，即
    \[
    \sigma
    =
    \sqrt{\frac{1}{D}\sum_{i=1}^{D}(x_i-\mu)^2}.
    \]

    \item $\operatorname{RMS}(x)$：以 $0$ 为参考点，衡量各个分量距离 $0$ 的整体大小，即
    \[
    \operatorname{RMS}(x)
    =
    \sqrt{\frac{1}{D}\sum_{i=1}^{D}(x_i-0)^2}
    =
    \sqrt{\frac{1}{D}\sum_{i=1}^{D}x_i^2}.
    \]
\end{itemize}

所以从公式形式上看，可以直观地把 RMS 理解成：把标准差中作为参考点的均值 $\mu$ 换成 $0$。

对最后一维长度为 $D$ 的向量 $x=[x_1,\ldots,x_D]$，RMSNorm 的完整公式是：
\[
\operatorname{RMSNorm}(x)_i
=w_i\frac{x_i}
{\sqrt{\frac{1}{D}\sum_{j=1}^{D}x_j^2+\varepsilon}},
\qquad i=1,\ldots,D.
\]
这里 $\varepsilon$ 是很小的正数，用来避免分母为 0 或过小时数值不稳定，常见形式是 $\sqrt{\operatorname{mean}(x^2)+\varepsilon}$。$w=[w_1,\ldots,w_D]\in\mathbb{R}^D$ 是参数向量，让模型在训练中自己决定每个维度的重要程度。所有通道共用同一个 RMS 分母，但各自有独立的 $w_i$；RMSNorm 没有 $\beta$ 偏移参数。

\section{手动算一遍 RMSNorm}

继续用 $x=[1,2,3,4]$。它的 RMS 为
\[
\operatorname{RMS}(x)
=
\sqrt{\frac{1^2+2^2+3^2+4^2}{4}}
=
\sqrt{7.5}
\approx2.739.
\]

因此归一化后的结果为
\[
\widehat{x}
=
\frac{x}{\operatorname{RMS}(x)}
=
\frac{[1,2,3,4]}{\sqrt{7.5}}
\approx
[0.365,0.730,1.095,1.461].
\]
在 $w=\mathbf{1}$、$\varepsilon=0$ 的简化条件下，四个坐标共同归一化后仍按 $1:2:3:4$ 分布；完整 RMSNorm 再由逐维权重调整尺度。


\section{两种归一化对比}

忽略参数的影响，令 $\varepsilon=0$，可训练缩放系数都设为 1，看同一个输入 $[1,2,3,4]$ 的结果：

\begin{center}
\small
\begin{tabular}{p{2.1cm}p{5.1cm}ccc}
\toprule
方法 & 输出 & 输出均值 & 输出 RMS & 居中\\
\midrule
原始输入 & $[1,2,3,4]$ & $2.500$ & $2.739$ & 否\\
LayerNorm & $[-1.342,-0.447,0.447,1.342]$ & $0$ & $1$ & 是\\
RMSNorm & $[0.365,0.730,1.095,1.461]$ & $0.913$ & $1$ & 否\\
\bottomrule
\end{tabular}
\end{center}

这个表只说明计算上的区别，不说明哪种方法一定更好。LayerNorm 同时控制均值和尺度；RMSNorm 只控制尺度。很多 decoder-only 模型用 RMSNorm 后训练效果没有下降，还省掉了一部分统计计算，所以慢慢流行起来。

\section{RMSNorm 在 Transformer 里的位置}

前面我们一直把 RMSNorm 当成一个单独的归一化操作来看，但在真正的 Transformer 中，它并不是独立存在的，而是夹在各个子层之前，负责先把输入向量的尺度整理稳定，再交给后面的计算模块。

现代 decoder-only Transformer 通常采用 \textbf{pre-norm} 结构。以一层 Transformer 为例，数据流大致是：

\[
X
\longrightarrow
\operatorname{RMSNorm}(X)
\longrightarrow
\operatorname{Attention}
\longrightarrow
\text{与原来的 }X\text{ 相加}.
\]

这里“与原来的 $X$ 相加”使用的是 Transformer 中常见的
\textbf{残差连接（residual connection）}。它的具体作用和原理后面会单独介绍，这里暂时只需要把它理解成：“我保留了一点大肠的味道，这样你才知道你吃的是大肠”。

写成公式就是：

\[
X_1
=
\underbrace{X}_{\text{原始 }X}
+
\underbrace{
\operatorname{Attention}
\bigl(\operatorname{RMSNorm}(X)\bigr)
}_{\text{经过 RMSNorm 和 Attention 后的新结果}}
\]
得到新的表示 $X_1$ 后，还会经过一次类似的过程：

\[
X_1
\longrightarrow
\operatorname{RMSNorm}(X_1)
\longrightarrow
\operatorname{MLP}
\longrightarrow
\text{与原来的 }X_1\text{ 相加},
\]

也就是

\[
X_2
=
X_1
+
\operatorname{MLP}
\bigl(\operatorname{RMSNorm}(X_1)\bigr).
\]
可以看到，只要在进入计算模块前，都会经过RMSNorm，所以无论来的数是小卡拉米还是大罗金仙，都会强行压缩成普通人，以确保每个计算组件的计算尺度都是类似的。
\end{document}
